\section{One-pass simulation of disturbing quantum instruments}
\label{sec:instruments65}
We now allow genuine state disturbance. This is a constructive
finite-horizon statement, not a finite-state classification. It is
separate from the classification by repeatable nondisturbing probes.

\subsection{Fixed physical data and operational error}
Fix $d\ge2$, a finite command set $\mathcal A$, finite outcome sets,
and an instrument $(\Phi_{a,y})_y$ for each $a\in\mathcal A$. Assume
a fixed Gaussian-rational Kraus description
\begin{equation}\label{eq:krausdata65}
 \Phi_{a,y}(X)=q^{-2}\sum_r M_{a,y,r}X M_{a,y,r}^*,\quad
 M_{a,y,r}\in\Z[i]^{d\times d},\quad
 \sum_{y,r}M_{a,y,r}^*M_{a,y,r}=q^2I
\end{equation}
with one fixed positive integer $q$. The initial density matrix is
$P_0/z_0$, where $P_0$ is Gaussian integral and positive, and
$z_0=\tr P_0>0$. Rationality here refers to this supplied finite
description; no assertion that every rational superoperator has a
Gaussian-rational Kraus decomposition is needed.

At each step an external classical controller chooses a command as
a function of the public history, possibly using its own randomness.
The instrument emits $y$ and updates the physical state accordingly.
The simulator receives the same kind of command and emits one outcome
immediately. It cannot reread consumed commands or emitted outcomes;
any information retained privately is charged. A common controller is
used when comparing the two experiments. Its memory is not simulator
memory, and no access to an unrecorded external history is granted to
the simulator.

For a complete control--outcome history $h$ at time $t$, let $A_h$
be the true subnormalized state, of trace equal to the probability
of $h$. Let $\widehat A_h$ be the analogous simulator matrix. Define
\begin{equation}\label{eq:cqerror65}
       \mathfrak D_t=\sum_h\|A_h-\widehat A_h\|_1.
\end{equation}
This is the trace norm of the difference of the block-diagonal
classical--quantum states, not an average conditioned on a rare
observed history. Histories of probability zero contribute the zero
matrix. Classical transcript total variation is at most
$\mathfrak D_t/2$. The same bound holds after any common,
history-dependent terminal measurement. The compared interface consists of the classical transcript
and its numerical matrix blocks. We do not claim total variation of
binary matrix encodings, nor simulation of an arbitrary external
entangled quantum memory.

\begin{theorem}[Charged simulation of rational instruments]\label{thm:instrumentstream65}
For every fixed input \eqref{eq:krausdata65} there is one randomized
classical one-pass program which, for all $N\ge2$ and $L\ge2$,
simulates every $N$-step adaptive control policy and satisfies
\begin{equation}\label{eq:instrumenterror65}
                  \sup_{\rm policies}\mathfrak D_t\le2^{-L}
                         \qquad(0\le t\le N).
\end{equation}
Its stored and returned numerical state is always in $\mathcal D_d$.
Writing
\[
 s_*=\min\{N,L+\lceil\log_2(N+1)\rceil\},
\]
its writable space is $O_{\rm input}(s_*)$, equivalently
$O(d^2s_*)$ with constants allowed to depend on the fixed Kraus data.
Command processing uses $O_{\rm input}(s_*^2)$ expected bit operations
and $O_{\rm input}(s_*)$ expected fair bits per step. Every step
terminates almost surely, with a worst-case space bound even during
rejection sampling. No exact-probability sampling oracle is used.

Exact simulation of the entire adaptive experiment, including its
conditional numerical matrix states, is possible in
$O_{\rm input}(N)$ writable bits. No positivity bound on outcome
probabilities, reversibility, physical return, or spectral gap is
assumed. For one-outcome instruments the program is deterministic.
\end{theorem}
The finite program and physical data are fixed before $N,L$ are
supplied. Parameter scanning is charged by its actual input length.
The theorem concerns a classical numerical simulation, not preparation
of a quantum state. It gives no uniform error bound for normalized
states conditional on arbitrary rare histories.

\subsection{An integer realization of every step}
For approximate simulation choose
\begin{equation}\label{eq:instrumentprecision65}
 b=L+\left\lceil\log_2\bigl(6d^2(N+1)\bigr)\right\rceil,\qquad
 B=2^b,\quad m=B+2d^2.
\end{equation}
Store $\mathcal R_B(P_0/z_0)$ as $P/m$ with an integral Hermitian
numerator $P$. Given command $a$, form
\begin{equation}\label{eq:instrumentintegers65}
 S_y=\sum_r M_{a,y,r}P M_{a,y,r}^*,\quad
 s_y=\tr S_y,\quad T=\sum_y s_y=q^2m.
\end{equation}
Emit $y$ with probability $s_y/T$. Only positive $s_y$ can be emitted.
Replace the state by $\mathcal R_B(S_y/s_y)$, using the divisions in
Lemma~\ref{lem:positiveround65}. Its denominator is again exactly $m$.

These rational probabilities have an explicit fair-bit implementation.
Let $\ell=\lceil\log_2T\rceil$, draw $\ell$ independent fair bits as an
integer $J\in\{0,\ldots,2^\ell-1\}$, and retry if $J\ge T$. If $T=1$
take $J=0$ without drawing a bit. Accept the unique outcome whose
cumulative interval of integer lengths $(s_y)_y$ contains $J$.
Acceptance probability is greater than $1/2$, except that for a power
of two it is one. Thus the expected number of bits is at most $2\ell$.
Discarded trials are not stored, and their count is not needed.
Zero-length intervals can never be selected.

There is also an exact mode. Start from $P_0$ and, on each command,
form $(S_y)_y$ by \eqref{eq:instrumentintegers65}, now with
$T=q^2\tr P$. Sample the same integer intervals and set $P\leftarrow S_y$.
The numerical state is $P/\tr P$. This update is precisely the
normalized physical branch, while the numerator itself is never
divided. After $t$ commands,
\begin{equation}\label{eq:exacttrace65}
                      1\le\tr P\le z_0 q^{2t}.
\end{equation}
It is an integer because all input matrices are Gaussian integral.
The upper bound follows from
$\tr S_y\le\sum_y\tr S_y=q^2\tr P$.
All entries have modulus at most $\tr P$ by positivity.

\subsection{Uniform control of the whole path}
\begin{proof}[Proof of Theorem~\ref{thm:instrumentstream65}]
In approximate mode, each stored state is legal by
Lemma~\ref{lem:positiveround65}, and every sampling row has total
mass one by \eqref{eq:krausdata65}. For an arbitrary common controller,
write $\pi(a\mid h)$ for its next command law. The true and approximate
subnormalized updates are
\[
 A_{hay}=\pi(a\mid h)\Phi_{a,y}(A_h),\qquad
 \widehat A_{hay}=\pi(a\mid h)
                \overline{\mathcal R}_B(\Phi_{a,y}(\widehat A_h)).
\]
This expression agrees with the actual emitted probabilities and
normalized successor matrices. For each history, the triangle
inequality and Lemma~\ref{lem:instrumentcontraction65} imply
\begin{align*}
 \sum_{a,y}\|A_{hay}-\widehat A_{hay}\|_1
 &\le\|A_h-\widehat A_h\|_1+
          \delta_B\tr\widehat A_h .
\end{align*}
Summing over histories and using $\sum_h\tr\widehat A_h=1$ gives
$\mathfrak D_{t+1}\le\mathfrak D_t+\delta_B$.
Initial rounding contributes $\mathfrak D_0\le\delta_B$.
Hence
\[
                  \mathfrak D_t\le(t+1)\delta_B
                    \le(N+1)6d^2/B\le2^{-L}.
\]
This proof works uniformly for every controller. A controller with
private randomness may be conditioned on its seed, followed by
integration; the displayed bound is independent of that seed.
Taking traces gives the classical total-variation assertion.
Appending a terminal instrument and applying
\eqref{eq:instrumentcontraction65} gives the assertion about terminal
measurements. No division by a small target probability has appeared.

In grid mode, the numerator $P$ and the denominator $m$ have
$O_{\rm input}(b+1)$ bits per entry. The fixed integral Kraus products,
the outcome traces, the cumulative intervals, and the integer $J$
have the same order. The numerators in the signed divisions have at
most $2b+O_{\rm input}(1)$ bits. A constant number, depending on the
fixed physical data, of matrix buffers suffices. Elementary long
division gives quadratic bit time and linear scratch space. Rejection
sampling uses at most $O_{\rm input}(b+1)$ space in every trial,
and has the stated expected bit and time bounds.

Exact mode reproduces every branch probability and conditional state
exactly. Equation~\eqref{eq:exacttrace65}, together with the fixed
Kraus coefficients, bounds all stored and temporary integers by
$O_{\rm input}(t+1)$ bits. Thus the exact simulation uses linear
space and has the stated conservative quadratic time bound per step.

Use grid mode when $b<N$ and exact mode otherwise. The choice is made
before allocating $B$: store $N$, scan $L$ with its value capped at
$N$, and compute the threshold with binary integer lengths. A
precision string of arbitrary length is consumed without being stored.
A counter verifying the number of commands uses $O(\log N)$ bits.
The fixed-radix numerical matrix output has
$O_{\rm input}(s_*)$ bits and requires no readable transcript tape.
These observations complete the charged-space proof.
\end{proof}

\begin{corollary}[Rational channels with legal matrix outputs]
\label{cor:channelstream65}
A fixed finite family of Gaussian-rational Kraus channels and a fixed
rational initial density matrix admit one deterministic one-pass
algorithm with pointwise trace-norm error at most $2^{-L}$ at every
prefix, legal outputs, and space
$O_{\rm input}(\min\{N,L+\log(N+1)\})$. Exact output uses
$O_{\rm input}(N)$ bits.
\end{corollary}
\begin{proof}
There is one outcome per command, so its probability is one and no
random bits are needed. The single block in
\eqref{eq:cqerror65} is the difference of the two numerical matrices.
Apply Theorem~\ref{thm:instrumentstream65}.
\end{proof}

The contraction and hybrid argument are classical. Their use here is
to give a fully specified positive integer simulator with bounded
space, exact rational emissions and no rare-outcome assumption.
This constructive theorem does not replace the repeatable-probe
hypothesis in the finite-response-quotient classification.
